% GENERATED FILE. Edit canonical lessons, then run npm run generate:books.
% canonical-source-hash: fnv1a64:3a4c52b826fd20a8
% canonical-lessons: GU-C519-sagdi, GU-C519-jhal, GU-C519-ac, GU-C519-varal, GU-C519-thaarvu

\chapter{Fire and Steam}
\label{ch:gu-fire-and-steam}
\hlchaptermodality{519}

\begin{chapteropening}
\textbf{By the end of this chapter:} \emph{I can talk about heat, flame and steam in the kitchen.}

Complete the last lesson of chapter 519: I can talk about heat, flame and steam in the kitchen.
\end{chapteropening}

\section[\texorpdfstring{\gu{સગડી} (sagḍī)}{સગડી (sagḍī)}]{\gu{સગડી} (sagḍī) — a charcoal stove}

\label{lesson:GU-C519-sagdi}



\begin{quote}
Before the new one: say the Gujarati for mould, then the Gujarati for to stick (something) on.
\end{quote}

\subsection*{You'll want to know: \gu{સગડી}}
\textbf{\gu{સગડી}} — \emph{sagḍī} — \textquotedblleft{}a charcoal stove\textquotedblright{}.

\begin{quote}
\gu{શિયાળામાં} \gu{દાદી} \gu{સગડી} \gu{પાસે} \gu{બેસતી} \gu{હતી}.
\end{quote}

\begin{quote}
\emph{śiyāḷāmā̃ dādī sagḍī pāse bestī hatī.}
\end{quote}

\begin{quote}
In winter my grandmother used to sit by the charcoal stove.
\end{quote}

A feminine noun.

\begin{cousinweb}
\textbf{\gu{સગડી}} (\emph{sagḍī}) No dictionary we use records an older form, so its history is not settled.
\end{cousinweb}

\begin{grammarlens}[title={What you've built}]
One more word for describing what you lived and felt.
\end{grammarlens}

\subsection*{Your turn}
\begin{itemize}
\raggedright
  \item \emph{Say it:} \emph{sagḍī}
  \item \emph{Say it:} \emph{sagḍī}, once more
  \item \emph{Say it:} \emph{śiyāḷāmā̃ dādī sagḍī pāse bestī hatī.} — In winter my grandmother used to sit by the charcoal stove.
  \item \emph{Recall:} say the Gujarati for mould, then the Gujarati for to stick (something) on, then say \emph{sagḍī} again
\end{itemize}

\begin{tcolorbox}[breakable,colback=teal!4,colframe=teal!35!black,title={Before you move on}]
What does \gu{સગડી} mean? (\textquotedblleft{}A charcoal stove\textquotedblright{}.) Say it once more.
\end{tcolorbox}
\section[\texorpdfstring{\gu{ઝાળ} (jhāḷ)}{ઝાળ (jhāḷ)}]{\gu{ઝાળ} (jhāḷ) — a tongue of flame}

\label{lesson:GU-C519-jhal}



\begin{quote}
Before the new one: say the Gujarati for to stick (something) on, then the Gujarati for a charcoal stove.
\end{quote}

\subsection*{You'll want to know: \gu{ઝાળ}}
\textbf{\gu{ઝાળ}} — \emph{jhāḷ} — \textquotedblleft{}a tongue of flame\textquotedblright{}.

\begin{quote}
\gu{ચૂલાની} \gu{ઝાળ} \gu{લાલ} \gu{અને} \gu{પીળી} \gu{છે}.
\end{quote}

\begin{quote}
\emph{cūlānī jhāḷ lāl ane pīḷī che.}
\end{quote}

\begin{quote}
The flame of the stove is red and yellow.
\end{quote}

A feminine noun.

\begin{cousinweb}
\textbf{\gu{ઝાળ}} (\emph{jhāḷ}) comes from \emph{*jhāla}, a form CDIAL reconstructs, and CDIAL lists the Gujarati word, \textquotedblleft{}flame, heat\textquotedblright{} (CDIAL 5379).
\end{cousinweb}

\begin{grammarlens}[title={What you've built}]
One more word for describing what you lived and felt.
\end{grammarlens}

\subsection*{Your turn}
\begin{itemize}
\raggedright
  \item \emph{Say it:} \emph{jhāḷ}
  \item \emph{Say it:} \emph{jhāḷ}, once more
  \item \emph{Say it:} \emph{cūlānī jhāḷ lāl ane pīḷī che.} — The flame of the stove is red and yellow.
  \item \emph{Recall:} say the Gujarati for to stick (something) on, then the Gujarati for a charcoal stove, then say \emph{jhāḷ} again
\end{itemize}

\begin{tcolorbox}[breakable,colback=teal!4,colframe=teal!35!black,title={Before you move on}]
What does \gu{ઝાળ} mean? (\textquotedblleft{}A tongue of flame\textquotedblright{}.) Say it once more.
\end{tcolorbox}
\section[\texorpdfstring{\gu{આંચ} (ā̃c)}{આંચ (ā̃c)}]{\gu{આંચ} (ā̃c) — cooking heat (the flame under a pot)}

\label{lesson:GU-C519-ac}



\begin{quote}
Before the new one: say the Gujarati for a charcoal stove, then the Gujarati for a tongue of flame.
\end{quote}

\subsection*{You'll want to know: \gu{આંચ}}
\textbf{\gu{આંચ}} — \emph{ā̃c} — \textquotedblleft{}cooking heat (the flame under a pot)\textquotedblright{}.

\begin{quote}
\gu{દૂધને} \gu{ધીમી} \gu{આંચ} \gu{પર} \gu{ઉકાળો}.
\end{quote}

\begin{quote}
\emph{dūdhne dhīmī ā̃c par ukāḷo.}
\end{quote}

\begin{quote}
Boil the milk on a low flame.
\end{quote}

A feminine noun.

\begin{cousinweb}
\textbf{\gu{આંચ}} (\emph{ā̃c}) is a tadbhava, inherited from Sanskrit \emph{arcís}, \textquotedblleft{}a flame, a ray\textquotedblright{}, and CDIAL lists the Gujarati word, \textquotedblleft{}heat\textquotedblright{} (CDIAL 635). A cook speaks of \textbf{\gu{ધીમી} \gu{આંચ}}, a low flame.
\end{cousinweb}

\begin{grammarlens}[title={What you've built}]
One more word for describing what you lived and felt.
\end{grammarlens}

\subsection*{Your turn}
\begin{itemize}
\raggedright
  \item \emph{Say it:} \emph{ā̃c}
  \item \emph{Say it:} \emph{ā̃c}, once more
  \item \emph{Say it:} \emph{dūdhne dhīmī ā̃c par ukāḷo.} — Boil the milk on a low flame.
  \item \emph{Recall:} say the Gujarati for a charcoal stove, then the Gujarati for a tongue of flame, then say \emph{ā̃c} again
\end{itemize}

\begin{tcolorbox}[breakable,colback=teal!4,colframe=teal!35!black,title={Before you move on}]
What does \gu{આંચ} mean? (\textquotedblleft{}Cooking heat (the flame under a pot)\textquotedblright{}.) Say it once more.
\end{tcolorbox}
\section[\texorpdfstring{\gu{વરાળ} (varāḷ)}{વરાળ (varāḷ)}]{\gu{વરાળ} (varāḷ) — steam}

\label{lesson:GU-C519-varal}



\begin{quote}
Before the new one: say the Gujarati for a tongue of flame, then the Gujarati for cooking heat.
\end{quote}

\subsection*{You'll want to know: \gu{વરાળ}}
\textbf{\gu{વરાળ}} — \emph{varāḷ} — \textquotedblleft{}steam\textquotedblright{}.

\begin{quote}
\gu{તપેલીમાંથી} \gu{ગરમ} \gu{વરાળ} \gu{નીકળે} \gu{છે}.
\end{quote}

\begin{quote}
\emph{tapelīmā̃thī garam varāḷ nīkḷe che.}
\end{quote}

\begin{quote}
Hot steam is coming out of the pot.
\end{quote}

A feminine noun.

\begin{cousinweb}
\textbf{\gu{વરાળ}} (\emph{varāḷ}) No dictionary we use records an older form, so its history is not settled.
\end{cousinweb}

\begin{grammarlens}[title={What you've built}]
One more word for describing what you lived and felt.
\end{grammarlens}

\subsection*{Your turn}
\begin{itemize}
\raggedright
  \item \emph{Say it:} \emph{varāḷ}
  \item \emph{Say it:} \emph{varāḷ}, once more
  \item \emph{Say it:} \emph{tapelīmā̃thī garam varāḷ nīkḷe che.} — Hot steam is coming out of the pot.
  \item \emph{Recall:} say the Gujarati for a tongue of flame, then the Gujarati for cooking heat, then say \emph{varāḷ} again
\end{itemize}

\begin{tcolorbox}[breakable,colback=teal!4,colframe=teal!35!black,title={Before you move on}]
What does \gu{વરાળ} mean? (\textquotedblleft{}Steam\textquotedblright{}.) Say it once more.
\end{tcolorbox}
\section[\texorpdfstring{\gu{ઠારવું} (ṭhārvũ)}{ઠારવું (ṭhārvũ)}]{\gu{ઠારવું} (ṭhārvũ) — to let cool}

\label{lesson:GU-C519-thaarvu}



\begin{quote}
Before the new one: say the Gujarati for cooking heat, then the Gujarati for steam.
\end{quote}

\subsection*{You'll want to know: \gu{ઠારવું}}
\textbf{\gu{ઠારવું}} — \emph{ṭhārvũ} — \textquotedblleft{}to let cool\textquotedblright{}.

\begin{quote}
\gu{મીરાએ} \gu{ગરમ} \gu{દૂધ} \gu{ઠાર્યું}.
\end{quote}

\begin{quote}
\emph{Mīrāe garam dūdh ṭhāryũ.}
\end{quote}

\begin{quote}
Mira let the hot milk cool.
\end{quote}

Past, with \gu{તેણે} (he or she …): \gu{તેણે} \gu{દૂધ} \gu{ઠાર્યું}, \emph{teṇe dūdh ṭhāryũ}.

\begin{cousinweb}
\textbf{\gu{ઠારવું}} (\emph{ṭhārvũ}) is the causative partner of \textbf{\gu{ઠરવું}}, \textquotedblleft{}to cool down\textquotedblright{}, and CDIAL lists it under \emph{*stharati}, \textquotedblleft{}is firm\textquotedblright{}, as \textquotedblleft{}to let cool\textquotedblright{} (CDIAL 13742).
\end{cousinweb}

\begin{grammarlens}[title={What you've built}]
One more word for describing what you lived and felt.
\end{grammarlens}

\subsection*{Your turn}
\begin{itemize}
\raggedright
  \item \emph{Say it:} \emph{ṭhārvũ}
  \item \emph{Say it:} \emph{ṭhārvũ}, once more
  \item \emph{Say it:} \emph{Mīrāe garam dūdh ṭhāryũ.} — Mira let the hot milk cool.
  \item \emph{Recall:} say the Gujarati for cooking heat, then the Gujarati for steam, then say \emph{ṭhārvũ} again
\end{itemize}

\begin{tcolorbox}[breakable,colback=teal!4,colframe=teal!35!black,title={Before you move on}]
What does \gu{ઠારવું} mean? (\textquotedblleft{}To let cool\textquotedblright{}.) Say it once more.
\end{tcolorbox}
